
\documentclass[a4paper]{article}
\usepackage{amsmath}
\usepackage{amssymb}

\begin{document}

\normalsize

\noindent \large Auteur: Dina Haajou \\
\noindent \large Studierichting: Informatica \\
\noindent \large Oefeningenreeks 6B nr. 1.3 \\

\medskip

\normalsize

Gegeven: $S_{14}$ van $\left( \dfrac{5}{8}, 1, \dfrac{11}{8}, \dfrac{7}{4}, \ldots \right)$ \\

De rij is rekenkundig want het verschil tussen opeenvolgende termen is \\ 

constant: \\

$1 - \dfrac{5}{8} = \dfrac{3}{8}$ \\

$\dfrac{11}{8} - 1 = \dfrac{3}{8}$ \\

$\Rightarrow v = \dfrac{3}{8}$ \\

De eerste term is $x_1 = \dfrac{5}{8}$ \\

De formule voor de partiële som van een rekenkundige rij is: \\

$S_n = \dfrac{n}{2} \cdot (2x_1 + (n - 1)v)$ \\

$S_{14} = \dfrac{14}{2} \cdot \left(2 \cdot \dfrac{5}{8} + (14 - 1) \cdot \dfrac{3}{8} \right)$ \\

$= 7 \cdot \left( \dfrac{10}{8} + \dfrac{39}{8} \right)$ \\

$= 7 \cdot \dfrac{49}{8} = \dfrac{343}{8}$

\begin{thebibliography}{999}
% \bibitem{}
\end{thebibliography}

\end{document}
